\uuid{F4IB}
\titre{Tangente}
\auteur{bodin ; hanani ; mzari}
\organisation{exo7}
\licence{CC BY-NC-SA 4.0}
\datecreate{2020-01-20}
\niveau{L1}
\module{Analyse}
\chapitre{Dérivabilité des fonctions réelles}
\sousChapitre{Autre}
\difficulte{2}
\theme{Dérivées, tangente}
\exo7qcmid{lille-1.derivees.013}

\contenu{
\begin{question}
\qclassification{124.00}
\qtags{motcle=tangente}
Soit $\displaystyle f(x)=x^4-3x^3+3x^2-x$. Quelles sont les bonnes réponses ?
\begin{answers}
    \good{Il existe $a\in ]0,1[$ tel que $f'(a)=0$.}
    \good{Il existe $a\in ]0,1[$ où la tangente à $\mathcal{C}_f$ en $a$ est une droite horizontale.}
    \bad{Il existe $a\in ]0,1[$ où la tangente à $\mathcal{C}_f$ en $a$ est une droite verticale.}
    \bad{$\mathcal{C}_f$ admet un point d'inflexion en $0$.}
\end{answers}
\begin{explanations}
La fonction $f$ est dérivable sur $\Rr$. En particulier, la tangente à $\mathcal{C}_f$ en un point $a\in \Rr$ ne peut être une droite verticale. Par ailleurs, $f(0)=f(1)=0$. Donc le théorème de Rolle implique l'existence de $a\in ]0,1[$ tel que $f'(a)=0$ et la tangente à $\mathcal{C}_f$ en ce point est une droite horizontale.
\end{explanations}
\end{question}
}
